\section{Pairs of large negative zeros}

For an integer $\ell\ge1$ write
\[
\rho_\ell=\frac{\ell+1}{\ell},\quad
a_\ell=\ell^{1-\ell},\quad
b_\ell=(\ell+1)^{1-\ell}+\ell^{-\ell},\quad
c_\ell=(\ell+1)^{-\ell},
\]
and set
\begin{equation}\label{eq:outer-separators}
\lambda_\ell^-=\frac{c_\ell}{2b_\ell},\qquad
\lambda_\ell^0=\frac{b_\ell}{2a_\ell},\qquad
\lambda_\ell^+=\frac{2b_\ell}{a_\ell},\qquad
X_\ell^\nu=\lambda_\ell^\nu\rho_\ell^n.
\end{equation}

\begin{lemma}\label{lem:outer}
Let $n\ge250$ and $L=\lfloor\sqrt{n/(4\log n)}\rfloor$. For $1\le\ell\le L$,
\[
\sgn P_n(-X_\ell^-)=\sgn P_n(-X_\ell^+)=(-1)^n,\qquad
\sgn P_n(-X_\ell^0)=(-1)^{n+1}.
\]
The intervals $(X_\ell^-,X_\ell^+)$ are disjoint. Consequently $P_n$ has at least $2L$ distinct negative zeros with moduli greater than $X_L^-$.
\end{lemma}
\begin{proof}
Fix $\ell$ and set $X=\lambda\rho_\ell^n$,
$Z=(\ell+1)^{2n}X^{-2\ell}$. The four ideal weights
$W_u(k,h)$ with $(k,h)\in\{\ell,\ell+1\}^2$ have signed sum and absolute sum
\[
ZQ(\lambda),\quad ZM(\lambda),\qquad
Q=a_\ell\lambda^2-b_\ell\lambda+c_\ell,\quad
M=a_\ell\lambda^2+b_\ell\lambda+c_\ell.
\]
Throughout $\lambda_\ell^-\le\lambda\le\lambda_\ell^+$,
\begin{equation}\label{eq:lambda-bounds}
\frac1{2(\ell+4)}<\lambda\le4.
\end{equation}
Indeed $b_\ell/c_\ell=\ell+1+\rho_\ell^\ell<\ell+4$, while
$b_\ell/a_\ell\le2$.

For fixed $h$, consecutive ideal $k$ ratios are
$W(k+1,h)/W(k,h)=\rho_k^n/(hX)$.
The outward right and left ratios beyond the two central columns are therefore at most
\[
q_R=2(\ell+4)e^{-n/(\ell+1)^2},
\qquad q_L=4n e^{-n/\ell^2},
\]
respectively; a left tail is absent at $\ell=1$.
The left estimate uses $h\le n$, sufficient for the actual triangular sum.
In either central column, the corresponding outward $h$ ratios are at most
$q_R$ and $16e^{-n/\ell^2}$.

Put $H=\sqrt{n/(4\log n)}\ge3$. Then $\ell+1\le4H/3$, $\ell+1\le\sqrt n/3$, and
\[
\frac n{(\ell+1)^2}\ge\frac94\log n,\qquad
\frac n{\ell^2}\ge4\log n,\qquad
q_R\le\frac76n^{-7/4}<\frac1{1000}.
\]
The other three ratios are also below $1/1000$. Summing the two geometric tails in each coordinate shows that the ideal mass outside the central four cells is at most
\[
\{(1+2/999)^2-1\}ZM<ZM/200.
\]
By \eqref{eq:bonferroni}, replacing the central ideal weights by the actual Eulerian products changes their sum by less than $ZM/1000$: each of the two relative errors is bounded by
$(n+1)e^{1-3\sqrt n}$, and $2(n+1)e^{1-3\sqrt n}<1/1000$.
Thus, after division by $(-X)^n Z$, the actual polynomial differs from $Q(\lambda)$ by less than $M(\lambda)/100$.

The three quadratic signs have uniform margins. Put
$u=(\ell+1)^{1-\ell}$, $v=\ell^{-\ell}$, $w=u/v$. Then
\[
\frac{uv}{a_\ell c_\ell}=\frac{\ell+1}{\ell},
\qquad w=(\ell+1)(\ell/(\ell+1))^\ell>\frac{\ell+1}{3}.
\]
For $\ell\ge2$, $w>1$, and
\[
\frac{b_\ell^2}{a_\ell c_\ell}
>\left(1+\frac1\ell\right)
 \left(\frac{\ell+1}{3}+\frac3{\ell+1}+2\right)
=\frac{(\ell+4)^2}{3\ell}\ge\frac{16}{3}.
\]
For $\ell=1$ the ratio is $8$. Substitution at the three values of $\lambda$ gives
\[
\frac{Q(\lambda^-)}{M(\lambda^-)}>\frac{32}{99},\qquad
-\frac{Q(\lambda^0)}{M(\lambda^0)}>\frac1{15},\qquad
\frac{Q(\lambda^+)}{M(\lambda^+)}>\frac13.
\]
For the middle bound, put $D=b_\ell^2/(4a_\ell)>4c_\ell/3$ and use
$(D-c_\ell)/(c_\ell+3D)>1/15$.
These margins prove the signs of the actual polynomial.
Finally
\[
\frac{X_{\ell+1}^+}{X_\ell^-}
<8(\ell+4)e^{-n/(\ell+1)^2}=4q_R<1.
\]
The two sign changes in each disjoint interval give the asserted roots.
\end{proof}

For the short intermediate range we can increase the cutoff.
\begin{lemma}\label{lem:finitebridge}
Let $250\le n\le1600$ and $L=\lceil\sqrt n/4\rceil$. The signs and disjoint intervals of Lemma~\ref{lem:outer} hold for every $\ell\le L$. Moreover, with $u=\log X_L^-$,
\begin{equation}\label{eq:finite-phase-match}
40<u<4\sqrt n,\qquad
(2L-n-1)\pi<\phi(-u)<(2L-n)\pi.
\end{equation}
\end{lemma}
\begin{proof}
Here $\ell\le10$ and
$n\ge n_\ell:=\max(250,16(\ell-1)^2+1)$.
The four outward ratios used above have the uniform bounds
\[
\begin{array}{c|cc}
 &\text{left}&\text{right}\\ \hline
k&.036&.015\\
h&.0006&.015
\end{array}
\]
and the central relative error is below $10^{-6}$.
These follow by evaluating the decreasing functions
$4n e^{-n/\ell^2}$, $2(\ell+4)e^{-n/(\ell+1)^2}$,
$16e^{-n/\ell^2}$ and $2(n+1)e^{-(n-\ell)/(\ell+1)}$
at the ten integer endpoints $n_\ell$. Exact rational bounds for these evaluations are part of Appendix~\ref{sec:verification}.
Each tail attaches to its own central row or column, so their common factors are
\[
\eta_k=\frac{.036}{1-.036},\qquad
\eta_h=\frac{.015}{1-.015},\qquad
(1+\eta_k)(1+\eta_h)-1<.054.
\]
Thus the quadratic error is less than $(.054+10^{-6})M<M/15$.
The same margins prove the three signs; disjointness follows from
$8(\ell+4)e^{-n/(\ell+1)^2}<.06$.
Since $\lambda_L^-<1$, $u<n/L\le4\sqrt n$.

We give the phase comparison in detail. For any $\ell\le L$, put
\[
D=2(\ell+1+\rho_\ell^\ell),\quad
u=n\log\rho_\ell-\log D,\quad
z=-u+i\pi,\quad \theta=\pi-\delta,\quad
\delta=\arctan(\pi/u),\quad y=\frac z{1+e^u}.
\]
The same ten endpoint checks give
\begin{equation}\label{eq:finite-phase-constants}
\begin{gathered}
u>40,\quad u>\frac{11n}{12(\ell+1)},\quad
u>\frac{n-\ell+3}{\ell+1},\\
\frac n{2(\ell+1)}>21,\qquad
\frac{(\ell+1)^2}{n}<\frac16,\\
(n+1)e^{-n/(\ell+2)}<10^{-6},\qquad e^{-u}<10^{-6}.
\end{gathered}
\end{equation}
For the first three inequalities one uses the lower bound
$u>n\log\rho_\ell-\log(2(\ell+4))$; all three differences increase with $n$.

Write $E(y)=\sum_{r=0}^{n-1}b_r y^r$,
$b_r=(n-r)!A(n,r+1)$. The first term to the left of $b_\ell y^\ell$ has ratio less than $p_0=11/20$. More generally, for every $1\le r\le\ell$,
\[
\frac{b_{r-1}}{b_r|y|}
\le\frac n{Du}\,\frac{1+e^{-u}}{1-10^{-6}}
       \left(\frac{\rho_\ell}{\rho_r}\right)^n
<p_0\left(\frac{\rho_\ell}{\rho_r}\right)^n.
\]
This uses $n-r+1\le n$ directly. Thus all further left ratios are less than
$p_0 e^{-n/\ell^2}<q_0=1/100$.
The first ratio to the right is at most
\[
\frac{2(\ell+4)(n/\ell+4)}{(n-\ell)(1-10^{-6})}
 e^{-n/(\ell+1)^2}<q_0.
\]
The rational prefactor and the exponential decrease with $n$; the endpoint check proves the bound. Since $E$ is the reversal of $F_n$, its coefficients are log-concave, so all later right ratios are no larger.
Consequently, for $J(y)=E(y)/(b_\ell y^\ell)$ on the entire semicircle of radius $|y|$,
\[
|J-1|<\frac{p_0+q_0}{1-q_0}<.57,\qquad
\sum_{r\ne\ell}|r-\ell|\frac{|b_r y^r|}{|b_\ell y^\ell|}
<\frac{p_0+q_0}{(1-q_0)^2}<.58.
\]
At angle $\theta=\pi-\delta$, $|\Im J|<.58\delta$, $\Re J>.43$, and hence
$|\arg J|<2\delta$. Control on the whole semicircle fixes the continuous argument without a multiple of $2\pi$. Reversal now gives
\[
\phi(-u)=-(n-\ell+1)\theta+\arg J.
\]
The upper phase bound follows from
$(n-\ell+3)\delta/\pi<(n-\ell+3)/u<\ell+1$.
For the lower bound, the trapezoid estimate for $1/t$ and $D>4$ yield
\[
n-\ell-1-\ell u
\ge\frac n{2(\ell+1)}-\ell-1+\ell\log D
>\frac n{2(\ell+1)}-1>20.
\]
On the other hand,
$(n-\ell-1)\pi^2/(3u^2)<48\pi^2(\ell+1)^2/(121n)<1$.
Using $\arctan t>t-t^3/3$ gives
$(n-\ell-1)\delta/\pi>\ell$. The two bounds imply
$(2\ell-n-1)\pi<\phi(-u)<(2\ell-n)\pi$.
Taking $\ell=L$ proves \eqref{eq:finite-phase-match}.
\end{proof}
